
\subsubsection{2.1 Multi-Dimensional LLM Trustworthiness Evaluation}

\textbf{TrustLLM} \citep{Sun2024TrustLLM} evaluates 16 decoder-only LLMs across six trustworthiness dimensions --- fairness (BBQ), robustness (ANLI, OOD), privacy, safety, and truthfulness --- using consistent evaluation protocols. TrustLLM provides the score matrix used in this study. It reports absolute per-model scores rather than asking whether rankings on one condition predict rankings on another.

\textbf{DecodingTrust} [Wang et al., 2023a] evaluates GPT-3.5 and GPT-4 on eight trustworthiness dimensions, finding that GPT-4 is more vulnerable to adversarial jailbreaking despite higher standard benchmark performance. This N = 2 rank reversal observation motivated the adversarial disruption hypothesis examined here, which the present N = 16 analysis falsifies. The pattern does not generalize across a diverse model population.

\textbf{HELM} \citep{Liang2022Holistic} provides holistic evaluation across 42 models and 7 metrics, establishing the multi-benchmark multi-model paradigm. Like TrustLLM and DecodingTrust, HELM reports within-scenario performance rather than cross-scenario predictive validity.

\subsubsection{2.2 Adversarial Robustness Benchmarks}

\textbf{AdvGLUE} \citep{Wang2021AdvGLUE} applies 14 adversarial attack methods to GLUE tasks, producing a benchmark explicitly designed to defeat models that pass in-distribution GLUE. This adversarial design philosophy underpinned the hypothesis that rank stability would be disrupted. The present results show that OOD robustness rank stability (partial $\rho$ = 0.868) is high at the model-ranking level, contradicting the design-level intuition.

\textbf{ANLI} \citep{Nie2020Adversarial} constructs adversarial NLI rounds (R1, R2, R3) through iterative human-and-model-in-the-loop processes, with each successive round designed to defeat models passing earlier rounds. The finding that partial $\rho_{ANLI}$ = 0.684 (p = 0.007) after MMLU control is positive indicates that adversarial construction disrupts performance levels but preserves relative model ordering.

\textbf{Wang et al.} [2023b] evaluate ChatGPT on AdvGLUE and ANLI against baseline models, finding ChatGPT advantages on OOD tasks but not computing cross-model rank correlations or controlling for general capability. The present study extends this to 16 models with capability control.

\subsubsection{2.3 Fairness Benchmark Design}

\textbf{BBQ} \citep{Parrish2022BBQ} constructs 50,000 QA items across nine social bias categories, with disambiguated contexts (where factual information resolves the answer) and ambiguous contexts (where only stereotype-consistent or stereotype-inconsistent guessing is possible). The disambiguated $\rightarrow$ ambiguous shift constitutes a natural in-distribution/OOD pair for fairness. The finding that they share near-perfect partial $\rho$ = 0.962 validates BBQ's cross-context construct consistency at the model-ranking level, while leaving open whether this reflects stable model-internal mechanisms or shared benchmark structure (see Section 6.2).

\subsubsection{2.4 Cross-Benchmark Predictive Validity}

\textbf{Gevers and Daelemans} [2026] provide the closest methodological precedent: rank correlations with leave-one-family-out cross-validation across 23 LLMs on commonsense benchmarks. This work is cited as a preprint that was independently assessed as plausible but has not been confirmed via library access at submission time. The present study extends the approach to trustworthiness dimensions, which differ from commonsense tasks in their safety stakes and the mechanistic hypothesis distinguishing latent bias from adversarial design.

\textbf{GLUE-X} \citep{Yang2023GLUEX} measures in-distribution to OOD accuracy gaps for encoder-only pretrained language models (PLMs), finding systematic performance degradation. GLUE-X does not compute cross-model rank correlations, and its model population --- encoder-only PLMs --- has zero overlap with the decoder-only LLMs in the present study. The GLUE $\rightarrow$ AdvGLUE benchmark pair was therefore structurally unavailable for this analysis; OOD robustness (TrustLLM Micro F1) was used as a proxy.

No prior study computes capability-controlled partial Spearman $\rho$ between in-distribution and OOD trustworthiness benchmark model rankings across multiple dimensions and 15 or more models.

